% 2-SAT Solver

Solves 2-SAT in $O(V + E)$ using Tarjan's SCC. Variables are 1-indexed. Use negative sign for negation (e.g., -u).

\begin{lstlisting}
struct TwoSat {
    int n, timer, scc_count;
    vector<vector<int>> adj;
    vector<int> tin, low, scc_id;
    vector<bool> in_stack, answer;
    stack<int> st;

    TwoSat(int n) : n(n), adj(2 * n + 1), tin(2 * n + 1, -1), low(2 * n + 1, -1),
                    scc_id(2 * n + 1, -1), in_stack(2 * n + 1, false), answer(n + 1, false),
                    timer(0), scc_count(0) {}

    void add_clause(int u, int v) {
        int not_u = (u > 0) ? u + n : -u;
        int act_u = (u > 0) ? u : -u + n;
        int not_v = (v > 0) ? v + n : -v;
        int act_v = (v > 0) ? v : -v + n;
        adj[not_u].push_back(act_v); adj[not_v].push_back(act_u);
    }

    void dfs(int u) {
        tin[u] = low[u] = ++timer;
        st.push(u); in_stack[u] = true;
        for (int v : adj[u]) {
            if (tin[v] == -1) {
                dfs(v); low[u] = min(low[u], low[v]);
            } else if (in_stack[v]) {
                low[u] = min(low[u], tin[v]);
            }
        }
        if (low[u] == tin[u]) {
            scc_count++;
            while (true) {
                int v = st.top(); st.pop(); in_stack[v] = false;
                scc_id[v] = scc_count;
                if (u == v) break;
            }
        }
    }

    bool solve() {
        for (int i = 1; i <= 2 * n; i++) if (tin[i] == -1) dfs(i);
        for (int i = 1; i <= n; i++) {
            if (scc_id[i] == scc_id[i + n]) return false;
            answer[i] = scc_id[i] < scc_id[i + n];
        }
        return true;
    }
};
\end{lstlisting}
